\section{Reachable fibers and the remaining pointwise endpoint}
\label{sec:v86-causal}
The original stopped transport reaches only part of the fiber at a
translated roof. Retaining that selection yields a smaller sufficient
optical charge than the unrestricted fractional majorant. No mass of
the original source is removed: all source trajectories are included
until their first true exit.

Let $\tau^-(y)$ be the backward lifetime in the same regular exact
event. Define the actual reachable-fiber drift
\begin{equation}\label{eq:v86-causal-drift}
 \mathcal V_h(t)=\int_0^h\int
 \one_{\{\tau^-(y)>s\}}\frac{I_m(y)}{z(y)}\,d\mu_{t+s}(y)\,ds.
\end{equation}
The integrand is defined on $z>0$; values on the source-null critical
set are immaterial. The Borel segment bijection makes its selection
exactly $a_j(\theta)<t<t+s<b_j(\theta)$. All external variables
use the common product kernels. Thus the formulas below hold outside
one null set in $(t,h)$ for each fixed $m,R,n,k$, and consequently
for almost every $t$ at each stated fixed width. Countably many
widths share a null set; no null set common to all radii is asserted.

\begin{theorem}[A source-level reachable-fiber comparison]
\label{thm:v86-causal-comparison}
For the original weighted exit and inverse-twist defects,
\begin{equation}\label{eq:v86-causal-source}
 E_h^\varepsilon(t)+O_h^\varepsilon(t)
                 \le C\mathcal F_h(t)+C\mathcal V_h(t).
\end{equation}
For the original short-segment volume density $q_\delta$ and outgoing
occupancy trace $\mathcal Z_\delta$,
\begin{equation}\label{eq:v86-causal-trace}
 |\mathcal Z_\delta(t)-q_\delta(t)|\le C\mathcal V_\delta(t).
\end{equation}
Moreover
\begin{equation}\label{eq:v86-causal-majorant}
 \mathcal V_h(t)\le C_\eta\mathsf P_{\eta,h}v^\eta(t).
\end{equation}
No equality between the two potentials is claimed.
\end{theorem}
\begin{proof}
On one maximal segment write $w(u)=\omega(\Psi_j(\theta,u))$.
On a survivor, $[w(t)-w(t+h)]_+$ is paid by the negative variation
up to $t+h$. At a grazing first exit $w(b-)=0$, so the same variation
pays the entire positive starting weight. At a nongrazing first exit,
$w(t)\le w(b-)+\int_t^b[-w']_+$. Divide by the initial angular
normalization $4\pi c_*\chi(t)$, use $W^\varepsilon\le1$ and
$\lambda_\circ\le\chi\le\Lambda_\circ$, and sum these disjoint
source cases. The bound $|X\omega|\le C\omega I_m/z$ and the
Borel map identify their complete variation with a submeasure of
\eqref{eq:v86-causal-drift}. The endpoint terms are exactly the
nongrazing first-exit assignment $\mathcal F_h$, up to the common
ratio of angular bounds. This proves \eqref{eq:v86-causal-source}.
A target $y$ occurs only when its segment contains the starting roof,
which is precisely $\tau^-(y)>s$; later re-entries do not occur.

For \eqref{eq:v86-causal-trace}, use the full angular density
$A_j=\omega/(4\pi c_*\chi)$ on each selected short segment.
Its established variation bound is
$|\partial_uA_j|\le CA_jI_m/z$. Integrate from $t$ to its actual
forward endpoint, whose distance is at most $\delta$. At every
intermediate roof its backward lifetime exceeds $s=u-t$.
The same Borel bijection gives multiplicity one. Tonelli and positive
exhaustion of compact regular subsegments prove both comparisons;
no inverse physical Jacobian has been replaced by one.

If $\tau^-(y)>s$, then $y=\Phi_sx$ for a point $x$ on the same
exact event. Radial growth gives $z(y)^2\ge2\lambda_\circ s$.
Thus $z(y)^{-1}\le(2\lambda_\circ)^{-\kappa}
 s^{-\kappa}z(y)^{-\eta}$. Discarding the reachable-fiber selector
only on the upper side proves \eqref{eq:v86-causal-majorant}.
It also supplies fixed-count local finiteness almost everywhere
from the retained fractional potential theorem.
\end{proof}

Split $\mathcal V_h=\mathcal V_h^{\rm mod}+\mathcal V_h^{\rm hi}$
using $I_m(y)\le s^{-\gamma}$ and its complement, with the same
fixed exponents as in Section~\ref{sec:v86-envelope}. Then
\begin{equation}\label{eq:v86-causal-hi}
 \mathcal V_h^{\rm hi}\le C_\eta\mathcal D_{\eta,\gamma,h}.
\end{equation}
The complete optical variation carried by this large-sensitivity
sector therefore obeys the finite-count uniform envelope already
proved; the source cases may include grazing first exits.
The estimate concerns their selected variation contribution, not
an unsupported claim that every such trajectory remains in that
sector for its entire lifetime.

\begin{theorem}[The smaller optical charge and its endpoint budget]
\label{thm:v86-endpoint}
The monotone limits
\[
 \mathfrak C_M=\lim_{h\downarrow0}\mathcal H_M(\mathcal V_h),\qquad
 \mathfrak C_M^{\rm mod}
       =\lim_{h\downarrow0}\mathcal H_M(\mathcal V_h^{\rm mod})
\]
exist in $[0,\infty]$ and satisfy
\begin{equation}\label{eq:v86-charge-comparison}
 \mathfrak C_M=\mathfrak C_M^{\rm mod}
       \le C_\eta\mathfrak A_{\eta,M},\qquad
 \mathfrak B_M\le\mathfrak Q_M+C\mathfrak C_M,\quad
 \mathfrak Q_M\le\mathfrak B_M+C\mathfrak C_M.
\end{equation}
For all sufficiently large fixed $B$ and sufficiently small fixed $r$,
\begin{align}
 \mathcal E_M\le{}&C_MB^{-1/192}
       +C\mathfrak B_M+C\mathfrak C_M^{\rm mod}
       +C_\eta r^{1/32}
       +C_\Theta(1+\varepsilon(B)^{-1})\sqrt r,\label{eq:v86-endpoint-B}\\
 \mathcal E_M\le{}&C_MB^{-1/192}
       +C\mathfrak Q_M+C\mathfrak C_M^{\rm mod}
       +C_\eta r^{1/32}
       +C_\Theta(1+\varepsilon(B)^{-1})\sqrt r.\label{eq:v86-endpoint-Q}
\end{align}
Constants in the second line may be enlarged. The conditions
$\mathfrak C_M^{\rm mod}=\mathfrak Q_M=0$ for every $M$ imply
the original full pointwise target. Neither condition is asserted
as proved by the large-sensitivity estimate.
\end{theorem}
\begin{proof}
Each potential is monotone in its upper integration limit. Positivity,
\eqref{eq:v86-causal-hi} and \eqref{eq:v86-large-height} give equality
of the two causal charges, including infinite values. The remaining
comparisons follow from \eqref{eq:v86-causal-majorant} and the two
positive inequalities in \eqref{eq:v86-causal-trace}, with the
complete tail inside $\mathcal H_M$ before taking its width to zero.

For $0<\delta<r$, radial growth and the fractional comparison give
\[
 \mathcal V_r-\mathcal V_\delta
       \le C_\eta\int_\delta^r s^{-\kappa}v^\eta(t+s)\,ds.
\]
Cover $[\delta,r]$ by finitely many dyadic annuli before the collision
limsup. The fixed-width input bounds their normalized sum by
$C_\eta r^{1-\kappa}$. Applying $\mathcal H_M$ first and then
letting $\delta\downarrow0$ proves
\begin{equation}\label{eq:v86-finite-head}
 \mathcal H_M(\mathcal V_r)
                 \le\mathfrak C_M+C_\eta r^{1-\kappa}.
\end{equation}
No infinite physical sum is moved through a limsup in this step.

Average \eqref{eq:v86-causal-source} in $0<h<r$ and use monotonicity
of $\mathcal V_h$. Insert the result into the complete source-shift
input and then the scalar positive remainder. The fixed-width bad
source contributes $O(B^{-1/192})$. The unweighted fixed-width
bound pays the unchanged guard and angular term. The established
surface sandwich gives
$\mathcal H_M(\overline{\mathcal F}_r)
\le\mathfrak B_M+C_\eta r^{1/32}$.
Equation~\eqref{eq:v86-finite-head} and $1-\kappa>1/32$ prove
\eqref{eq:v86-endpoint-B}; \eqref{eq:v86-charge-comparison} proves
\eqref{eq:v86-endpoint-Q}. Under the stated zero-charge conditions,
choose $r(B)=c\varepsilon(B)^\Gamma$ with $\Gamma>2$.
For each $B$ the width and all auxiliary bands are fixed before the
collision limit; auxiliary bands are enlarged afterwards and $B$
is taken last. The original target follows, including its arithmetic
zero classes without normalization.
\end{proof}

The new finite-count theorem excludes the entire large-sensitivity
microscopic sector by a uniform essential-height envelope, rather
than by fixed-scale limiting bounds. What remains is the reachable
moderate-sensitivity optical concentration and the original
zero-residence volume concentration. The reachability constraint is
physical, not a chosen collar. Proving these last two quantities zero
still requires a source-specific estimate beyond those proved here.
